\begin{graphicspathcontext}{{./chapters/ciad/examples/imgs/},{./chapters/ciad/examples/imgs/auto/},\old}

\sidecite{Dani2024}
\begin{frame}[t]{Objectifs du réseau de capteurs}
	\begin{center}
		\titlebox*[.3\linewidth]{Multi-capteurs}{Comment intégrer plusieurs capteurs et plusieurs types de capteurs dans le système ?}
		\titlebox*[.3\linewidth]{Multi-méthodes}{Comment prendre en charge différents types de méthodes d'analyse d'images dans le système ?}
		\titlebox*[.3\linewidth]{Passage à l'échelle}{Holonique multi-agents, conception d'un système scalable de traitement et d'apprentissage multi-méthodes et multi-caméras} \\[.25cm]
		\includegraphics[width=.7\linewidth]{hsn-sn1} \\
		{\tiny Manjah, D. (2024). Stream Based Active Distillation and Holonic Architecture for Scalable Camera Systems. PhD thesis, UCLouvain, Belgium}
	\end{center}
\end{frame}

\sidecite{Dani2024}
\begin{frame}[t]{Organisations dans le réseau de capteurs}
	\vspace{-.17cm}
	\begin{columns}
		\begin{column}{.4\linewidth}
			\smaller\smaller
			\begin{block}{Plateforme cyber-physique - CPP}
				\emph{Fournir des perceptions et gérer la qualité et les ressources}
				\begin{description}
					\item[Observateur] délivrer des perceptions grâce aux données du capteur
					\item[Capteur] caméra ou système cyber-physique interne
					\item[Fournisseur de ressources] répartition équitable des ressources
				\end{description}
			\end{block}
			\begin{block}{Apprentissage}
				\emph{Améliorer le comportement du capteur avec l'apprentissage par distillation active}
				\begin{description}
					\item[Étudiant] algorithme du capteur pour une zone ou méthode spécifique
					\item[Enseignant] modèle général pour superviser les étudiants
				\end{description}
			\end{block}
		\end{column}
		\begin{column}{.6\linewidth}
			\raggedleft
			\includegraphics[width=\linewidth]{hsn-sn2} \\[.5cm]
			\includegraphics[width=.9\linewidth]{hsn-sn3}
		\end{column}
	\end{columns}
\end{frame}

\sidecite{Dani2024}
\figureslide{Mélange des organisations de service et d'apprentissage dans l'holarchie}{hsn-sn4}

\sidecite{Dani2024}
\begin{frame}{Résultats expérimentaux}
	\begin{block}{Expériences}
		\begin{description}
			\item 40 caméras (jeu de données AI-City Tracking Challenge)
			\item[Mono-méthode] YOLO v4 entraîné sur le jeu de données MIO-TCD
			\item[Multi-méthode] YOLO v4 + soustraction de fond
		\end{description}
	\end{block}
	\vspace{.5cm}
	\begin{tabularx}{\linewidth}{X|X|X|X|X}
		\hline
		\textbf{Méthode} & \textbf{HOTA} & \textbf{Méthode (sec)} & \textbf{Agent (sec)} & \textbf{Total (sec)} \\
		\hline\hline
		DNN uniquement & $0.36$ & $2327 \pm 21$ & $10 \pm 1$ & $2337$ \\
		\hline
		Multi-méthode & $0.36$ & $2169 \pm 31$ & $12 \pm 1$ & $2182$ \\
		\hline\hline
		Différence & $0$ & $-158$ & $2$ & $-155$ \\
		\hline
	\end{tabularx}
	\begin{center}
		{\small 10 simulations, Intel(R) Core(TM) i7-8700K CPU 3.70GHz, GPU NVIDIA GeForce GTX 1080}
	\end{center}
\end{frame}

\end{graphicspathcontext}

\endinput
